% ============================================================
% ch26.tex —— 第 26 章 球面上的直线：大圆与测地线
% ============================================================
\chapter{球面上的直线：大圆与测地线}

初中几何有一条铁律盖了你三年的书桌：\textbf{三角形内角和等于 $180^\circ$}。这一章我们把它连锅端——不是推翻，而是\textbf{划定边界}：它在哪种世界里成立，在哪种世界里崩塌。崩塌的方式本身，就是新几何的出生证明。而这一切的起点，是一个天真的问题：\textbf{球面上的"直线"是什么？}

\section{球面上没有"直"，只好发明}

\begin{kaiwei}
从北京飞纽约，为什么航线画出来\textbf{绕过北极圈}，而不是沿着纬线横跨太平洋？"两点之间线段最短"——可你在球面上，画不了直线。
\end{kaiwei}

平面上的直线有一个身份证明：\textbf{它是两点间的最短路径}（也是"唯一不拐弯的路"）。球面上没有直尺，但"最短"这个身份可以\textbf{绕道继承}：拿一根橡皮筋，两端钉在球面两点，任其绷紧收缩——它停在哪里，哪里就是球面上的"直线"。动手实验（地球仪 + 细线，两分钟）：

\begin{itemize}
\item 钉北京、纽约：橡皮筋\textbf{拱过北极附近}——不是沿纬线走太平洋！这就是航线的答案。
\item 钉任意两点：橡皮筋永远停在某条"\textbf{最大圆}"（球心也在这个圆的圆心上的圆）的弧上。
\item 反面试验：沿北纬 $40^\circ$ 纬线绷线——绷不住，橡皮筋会滑向大圆。纬线（除赤道外）\textbf{不是}球面直线。
\end{itemize}

\begin{faming}
\textbf{土名}：绷紧的线、球面上的直线 $\to$ \textbf{正式名}：\textbf{测地线}（geodesic，希腊语"划分大地"——大地测量的直线）。球面上的测地线恰好是\textbf{大圆}（great circle）：过球心的平面截球面所得的圆。赤道、所有经线圈都是大圆；纬线除赤道外都不是（它们围的圆心不在球心）。
\end{faming}

北京（北纬 $40^\circ$、东经 $116^\circ$）到纽约（北纬 $41^\circ$、西经 $74^\circ$）：大圆航线约 \textbf{11\,000} 公里，而沿北纬 $40^\circ$ 纬线走要约 $16\,200$ 公里（经度差 $190^\circ\times111\,\text{km}\times\cos40^\circ$）——\textbf{差五千多公里}。航空公司用真金白银投票：跨太平洋航班全部"走高纬"。为什么纬线是弯路？沿北纬 $40^\circ$ 走的旅人\textbf{一直在向左微拐}（相对测地线而言）——每一步都在偏离"绷紧的方向"，步步吃亏。

\section{球面三角形：内角和的叛变}

现在用大圆当"直边"，在球面上画三角形。选一个最好算的：\textbf{顶点取北极 $N$、赤道上经度相差 $90^\circ$ 的两点 $A$、$B$}。三条边：$NA$ 是经线（大圆 \checkmark），$NB$ 是经线（大圆 \checkmark），$AB$ 是赤道弧（大圆 \checkmark）——名门纯正的"球面直线三角形"。量三个角（角 $=$ 两边切线方向的夹角）：

\begin{itemize}
\item $A$ 处：经线与赤道\textbf{垂直}（经线指向正北、赤道指向正东）——$90^\circ$；
\item $B$ 处：同样 $90^\circ$；
\item $N$ 处：两条经线的夹角就是它们的\textbf{经度差} $90^\circ$（在北极上空俯视：经线像钟面辐条，差 $90^\circ$ 经度就是差 $90^\circ$ 角）。
\end{itemize}

\[
\boxed{\;\text{内角和} = 90^\circ + 90^\circ + 90^\circ = 270^\circ \;}
\]

\begin{tikzpicture}[scale=1.35]
  \draw (0,0) circle (1.5);
  \draw[dashed] (-1.5,0) arc (180:360:1.5 and 0.45);
  \draw (-1.5,0) arc (180:0:1.5 and 0.45);
  \fill (0,1.5) circle (1.3pt) node[above] {$N$};
  \fill (1.06,-0.32) circle (1.3pt) node[right] {$A$};
  \fill (-1.06,-0.32) circle (1.3pt) node[left] {$B$};
  \draw[thick,shenlan] (0,1.5) -- (1.06,-0.32) -- (-1.06,-0.32) -- cycle;
  \node[shenlan] at (0.15,-0.75) {\scriptsize 三条边都是大圆（经线$\times$2 + 赤道）};
  \node at (0,-2.15) {\scriptsize 内角和 $= 90^\circ+90^\circ+90^\circ = 270^\circ$};
\end{tikzpicture}

"$180^\circ$"当场叛变。而且叛变有\textbf{规律}——不是乱叛。再画一个"半边天大三角"：顶点 $N$、赤道上经度差 $180^\circ$ 的两点（三条边：两条相对的经线 + 半条赤道）。三个角：$A$ 处 $90^\circ$、$B$ 处 $90^\circ$、$N$ 处 $180^\circ$（两条相对经线在北极成一条直线）——内角和 $360^\circ$。三角更胖，超标更多。量出规律（吉拉德，1629 年）：

\begin{dingli}[吉拉德定理：角盈正比于面积]
球面三角形的内角和超出 $180^\circ$ 的部分（\textbf{角盈} $E$）与面积 $A$ 成正比，比例系数是半径平方的倒数：
\[
E = \frac{A}{R^2} \qquad (E \text{ 以弧度计}) .
\]
\end{dingli}

两案验算。案一（刚才的 $270^\circ$ 三角形）：角盈 $90^\circ = \tfrac\pi2$ 弧度，反推面积 $= R^2\cdot\tfrac\pi2$。实际盘面积：这个三角形恰好是\textbf{八分之一个球面}（三条互相垂直的半径切出的"卦限"）：$\tfrac18\times4\pi R^2 = \tfrac{\pi R^2}2$ \checkmark。案二（半边天大三角）：角盈 $180^\circ = \pi$，面积应为 $\pi R^2$；实际它由两条相对的经线与半条赤道围成，占北半球的一半——四分之一个球面，面积恰为 $\pi R^2$ \checkmark。全局再对账：球面被八个"卦限三角形"瓜分，八块角盈合计 $8\times\tfrac\pi2 = 4\pi = \tfrac{4\pi R^2}{R^2} = \tfrac{\text{总面积}}{R^2}$——\textbf{整个球面的角盈预算，分毫不差}。

\section{为什么你在地球上活了两千年没发现}

$180^\circ$ 从没在你眼前失灵过，原因是一个除法：角盈 $= \dfrac{\text{面积}}{R^2}$，而\textbf{地球的 $R$ 太大了}。算一笔操场账：操场上的三角形面积顶多 $10^4$ 平方米，$R^2 = (6.371\times10^6)^2 \approx 4.05\times10^{13}$ 平方米：
\[
E \approx \frac{10^4}{4\times10^{13}} = 2.5\times10^{-10}\ \text{弧度} \approx 0.000000014^\circ .
\]
任何量角器都会告诉你"就是 $180^\circ$"。\textbf{不是欧几里得错了，是你站的地方太小}——角盈按面积批发，摊到每个局部就稀薄到比仪器精度还低。这不是狡辩，是可以反过来的科学方法：\textbf{测角盈，就是在测你所在球面的半径}。高斯真的主持过一场大规模测地三角测量（汉诺威王国，1820 年代），传闻他暗中检验过"角盈是否为零"——用山顶三角形给宇宙做体检的第一次尝试（真假存疑，但思想是他的）。

这一课的价值超出几何：\textbf{局部经验无法裁决全局法则}。桌面上的几何实验全部支持欧几里得，可桌面太小。人类花了两千年才从"桌面几何"里抬起头——第 28 章把这堂课补完（宇宙的形状也要靠"测角盈"裁决）。

\section{给篇五装上导航}

本章收获列成装备清单，供后两章取用：

\begin{enumerate}
\item \textbf{弯曲世界里的"直线"$=$ 最短线 $=$ 测地线}（绷紧的橡皮筋）：平面 $=$ 直线，球面 $=$ 大圆；
\item \textbf{内角和是弯曲的探测器}：$> 180^\circ$ 说明世界正弯（球面型，第 27 章"正曲率"）；$< 180^\circ$ 预示反弯（马鞍型，"负曲率"）；$= 180^\circ$ 才是平面；
\item \textbf{角盈 $\propto$ 面积 $/ R^2$}：弯曲是"按面积批发"的，摊到局部就稀薄——这解释了"小地方测不出弯曲"，也提供了"大地方测弯曲"的钥匙。
\end{enumerate}

\begin{dongshou}
\begin{itemize}
  \yisi 三个顶点取坐标轴方向的单位点 $(1,0,0), (0,1,0), (0,0,1)$ 的球面三角形（单位球面）：证明三条边两两垂直（每个角 $90^\circ$，内角和 $270^\circ$），并用吉拉德公式核对面积：角盈 $\tfrac\pi2$，面积 $\tfrac\pi2 = \tfrac18\times4\pi$ \checkmark——恰好八分之一个球面。
  \yier 用地球仪和细线验证三个"大圆航线"：北京--莫斯科、东京--洛杉矶（惊喜：拱上白令海附近的阿留申群岛）、悉尼--圣地亚哥（拱向南极海域）。每个都先猜后测，记下橡皮筋的实际走向与你的直觉差多少（"南半球航线向南拱"是固定节目）。
  \yier 证明球面上\textbf{不存在平行线}：任何两条大圆必相交。（提示：大圆 $=$ 过球心的平面截球面；两个过球心的平面必相交于一条过球心的直线；这条直线与球面交于两点——两"直线"恰交于一对对径点。具体例：任何两条经线在北极、南极相遇。）
  \yier 操场账再算一档：一座城市尺度的三角形（边长 $10$ km，面积约 $43$ km$^2$）的角盈是多少弧度？($E \approx 4.3\times10^7/4.05\times10^{13} \approx 1.06\times10^{-6}$ 弧度 $\approx 0.22$ 角秒——现代精密测角能测到 $0.1$ 角秒，\textbf{城市尺度的测地已经能看见地球的弯曲}，这不是思想实验，是大地测量学的日常。）
\end{itemize}
\end{dongshou}

\begin{shaonao}
\textbf{球面圆周实验（下一章的主角预演）。}球面上的"圆"定义：到定点（圆心）的\textbf{测地距离}相等的点集。以北极 $P$ 为心、测地半径 $\rho$（沿经线量的弧长，弧度计）画圆——它就是纬度 $\left(90^\circ - \rho\right)$ 的纬线。算它的周长 $C$（三维空间里量）与测地半径 $\rho$ 的关系：

(1) 纬线的三维半径（到地轴的距离）$= R\sin\rho$（画出剖面图，直角三角形斜边 $R$、角 $\rho$ 在球心）。

(2) 周长 $C = 2\pi R\sin\rho$；而"圆周长公式期望值"$= 2\pi\times(\text{测地半径}) = 2\pi R\rho$。

(3) 比较：因 $\sin\rho < \rho$（$0 < \rho$ 时，正弦落在对角线下方），得
\[
\boxed{C < 2\pi\times(\text{测地半径})} \qquad (\text{球面上的圆"缩水"}) .
\]
代入 $\rho = \tfrac\pi4$（$45^\circ$）：$C = 2\pi R\times0.7071$，期望 $2\pi R\times0.7854$，\textbf{缩水 $10\%$}。一只贴着球面爬行的蚂蚁量出这个缺口，不必起飞就知道世界是弯的——下一章，它靠这个缺口测出弯曲的\textbf{精确数值}。
\end{shaonao}

\begin{chengshi}
吉拉德定理我们只做了验算与对账，没给证明（证明需要球面三角的面积公式或高斯-博内定理的特例，高中竞赛上层内容）；测地线的严格定义（变分法：长度泛函的极值曲线）放行大学；"绷紧的线 $=$ 大圆"在球面上有对称性论证（橡皮筋最短路径的旋转对称性逼出大圆），一般曲面则需求解测地线微分方程。另外"飞机大圆航线"假设地球是完美球体，实际是扁椭球，极点方向更扁、航线计算用更精细的模型修正千分之几。
\end{chengshi}

球面正弯、平面零弯、内角和超标与面积挂钩。但"弯"本身到底是什么？能不能\textbf{不离开曲面}就测出弯曲？下一章，一只不肯起飞的蚂蚁，替全人类回答这个问题。
